% GENERATED FILE. Edit canonical lessons, then run npm run generate:books.
% canonical-source-hash: fnv1a64:d3cff9b3fbf592b1
% canonical-lessons: TA-C177-vilaivu, TA-C177-piraccinai, TA-C177-tirvu, TA-C177-nokkam, TA-C177-anupavam

\chapter{Result, Problem, Solution, Purpose, Experience}
\label{ch:ta-result-problem-solution-purpose-experience}
\hlchaptermodality{177}

\begin{chapteropening}
\textbf{By the end of this chapter:} \emph{I can say a result, a problem, a solution, a purpose and an experience.}

Complete the last lesson of chapter 177: I can say a result, a problem, a solution, a purpose and an experience.
\end{chapteropening}

\section[\texorpdfstring{\ta{விளைவு} (viḷaivu)}{விளைவு (viḷaivu)}]{\ta{விளைவு} (viḷaivu) — a result, a consequence}

\label{lesson:TA-C177-vilaivu}



\begin{quote}
Before the new one: say the Tamil for for example, then the Tamil for otherwise.
\end{quote}

\subsection*{You'll want to know: \ta{விளைவு}}
\textbf{\ta{விளைவு}} — \emph{viḷaivu} — \textquotedblleft{}a result, a consequence\textquotedblright{}.

\begin{grammarlens}[title={What you've built}]
One more word for this chapter.
\end{grammarlens}

\subsection*{Your turn}
\begin{itemize}
\raggedright
  \item \emph{Say it:} \emph{viḷaivu}
  \item \emph{Say it:} \emph{viḷaivu}, once more
  \item \emph{Say it:} say \emph{viḷaivu}
  \item \emph{Recall:} say the Tamil for for example, then the Tamil for otherwise, then say \emph{viḷaivu} again
\end{itemize}

\begin{tcolorbox}[breakable,colback=teal!4,colframe=teal!35!black,title={Before you move on}]
What does \ta{விளைவு} mean? (\textquotedblleft{}A result, a consequence\textquotedblright{}.) Say it once more.
\end{tcolorbox}
\section[\texorpdfstring{\ta{பிரச்சினை} (piracciṉai)}{பிரச்சினை (piracciṉai)}]{\ta{பிரச்சினை} (piracciṉai) — a problem}

\label{lesson:TA-C177-piraccinai}



\begin{quote}
Before the new one: say the Tamil for otherwise, then the Tamil for a result.
\end{quote}

\subsection*{You'll want to know: \ta{பிரச்சினை}}
\textbf{\ta{பிரச்சினை}} — \emph{piracciṉai} — \textquotedblleft{}a problem\textquotedblright{}.

\begin{grammarlens}[title={What you've built}]
One more word for this chapter.
\end{grammarlens}

\subsection*{Your turn}
\begin{itemize}
\raggedright
  \item \emph{Say it:} \emph{piracciṉai}
  \item \emph{Say it:} \emph{piracciṉai}, once more
  \item \emph{Say it:} say \emph{piracciṉai}
  \item \emph{Recall:} say the Tamil for otherwise, then the Tamil for a result, then say \emph{piracciṉai} again
\end{itemize}

\begin{tcolorbox}[breakable,colback=teal!4,colframe=teal!35!black,title={Before you move on}]
What does \ta{பிரச்சினை} mean? (\textquotedblleft{}A problem\textquotedblright{}.) Say it once more.
\end{tcolorbox}
\section[\texorpdfstring{\ta{தீர்வு} (tīrvu)}{தீர்வு (tīrvu)}]{\ta{தீர்வு} (tīrvu) — a solution}

\label{lesson:TA-C177-tirvu}



\begin{quote}
Before the new one: say the Tamil for a result, then the Tamil for a problem.
\end{quote}

\subsection*{You'll want to know: \ta{தீர்வு}}
\textbf{\ta{தீர்வு}} — \emph{tīrvu} — \textquotedblleft{}a solution\textquotedblright{}.

\begin{grammarlens}[title={What you've built}]
One more word for this chapter.
\end{grammarlens}

\subsection*{Your turn}
\begin{itemize}
\raggedright
  \item \emph{Say it:} \emph{tīrvu}
  \item \emph{Say it:} \emph{tīrvu}, once more
  \item \emph{Say it:} say \emph{tīrvu}
  \item \emph{Recall:} say the Tamil for a result, then the Tamil for a problem, then say \emph{tīrvu} again
\end{itemize}

\begin{tcolorbox}[breakable,colback=teal!4,colframe=teal!35!black,title={Before you move on}]
What does \ta{தீர்வு} mean? (\textquotedblleft{}A solution\textquotedblright{}.) Say it once more.
\end{tcolorbox}
\section[\texorpdfstring{\ta{நோக்கம்} (nōkkam)}{நோக்கம் (nōkkam)}]{\ta{நோக்கம்} (nōkkam) — a purpose, an aim}

\label{lesson:TA-C177-nokkam}



\begin{quote}
Before the new one: say the Tamil for a problem, then the Tamil for a solution.
\end{quote}

\subsection*{You'll want to know: \ta{நோக்கம்}}
\textbf{\ta{நோக்கம்}} — \emph{nōkkam} — \textquotedblleft{}a purpose, an aim\textquotedblright{}.

\begin{grammarlens}[title={What you've built}]
One more word for this chapter.
\end{grammarlens}

\subsection*{Your turn}
\begin{itemize}
\raggedright
  \item \emph{Say it:} \emph{nōkkam}
  \item \emph{Say it:} \emph{nōkkam}, once more
  \item \emph{Say it:} say \emph{nōkkam}
  \item \emph{Recall:} say the Tamil for a problem, then the Tamil for a solution, then say \emph{nōkkam} again
\end{itemize}

\begin{tcolorbox}[breakable,colback=teal!4,colframe=teal!35!black,title={Before you move on}]
What does \ta{நோக்கம்} mean? (\textquotedblleft{}A purpose, an aim\textquotedblright{}.) Say it once more.
\end{tcolorbox}
\section[\texorpdfstring{\ta{அனுபவம்} (aṉupavam)}{அனுபவம் (aṉupavam)}]{\ta{அனுபவம்} (aṉupavam) — an experience}

\label{lesson:TA-C177-anupavam}



\begin{quote}
Before the new one: say the Tamil for a solution, then the Tamil for a purpose.
\end{quote}

\subsection*{You'll want to know: \ta{அனுபவம்}}
\textbf{\ta{அனுபவம்}} — \emph{aṉupavam} — \textquotedblleft{}an experience\textquotedblright{}.

\begin{grammarlens}[title={What you've built}]
One more word for this chapter.
\end{grammarlens}

\subsection*{Your turn}
\begin{itemize}
\raggedright
  \item \emph{Say it:} \emph{aṉupavam}
  \item \emph{Say it:} \emph{aṉupavam}, once more
  \item \emph{Say it:} say \emph{aṉupavam}
  \item \emph{Recall:} say the Tamil for a solution, then the Tamil for a purpose, then say \emph{aṉupavam} again
\end{itemize}

\begin{tcolorbox}[breakable,colback=teal!4,colframe=teal!35!black,title={Before you move on}]
What does \ta{அனுபவம்} mean? (\textquotedblleft{}An experience\textquotedblright{}.) Say it once more.
\end{tcolorbox}
